\chapter*{\centering{\large{ABSTRACT}}}
\singlespacing{}

\textbf{Prabowo Darmawi}, Fish Length Measurement Using the \emph{Harris-Corner Detection} Method and Fish Weight Estimation Using a Derivative Method of the \emph{Length-Weight Relationship} Based on Digital Image Processing. 
Computer Science Study Program, Faculty of Mathematics and Natural Sciences, Universitas Negeri Jakarta, May 2026.
\\
\\
As technology continues to develop, the aquaculture sector in Indonesia still faces
problems in measuring fish length and weight, which are commonly performed
manually. In this study, the author applies the \emph{Harris-Corner Detection}
method to measure fish length and a simple derivation of the
\emph{Length-Weight Relationship} method to estimate fish weight based on digital
image processing. The first step in measuring fish length and estimating fish weight is inputting the
fish image. Furthermore, the \emph{Harris-Corner Detection} method is used to
detect corner points in the fish image, followed by the use of a
\emph{Bounding-Box} method to select corner points and calculate fish length.
Finally, a simple derivation of the \emph{Length-Weight Relationship} is used to
estimate fish weight. The result of this study is an algorithm capable of calculating fish length and
weight based on image input.
\\
\\
\textbf{Keywords}: Corner Detection, Harris-Corner Detection, Bounding-Box, Length-Weight Relationship, Fish.